% ---------------------------------------------------------------- Preface
\unnumberedchapter{Preface}
\chaptersubtitle{How this book is organized, and how to read it}
\chapterrule

Building software is a lifecycle, and each part of it is easiest to understand when it sits in its
proper place. This book follows that lifecycle from end to end: from the moment an idea appears in
someone's head to the moment it runs in production. It concentrates on \emph{the craft of building
software}. It shows where teams, roles and their rituals fit into the process, but how to run a
team, and how to navigate a career, are left to other books.

\section*{The shape of the book}

The book opens with a map and closes by returning to it. \textbf{How Software Gets Built}, the
opening chapter, lays out the whole process before any of its parts: the loop every piece of
software goes around, the pieces a running application is made of, and the people who do each part.
Between the map and the closing chapter are four parts, one for each phase of the lifecycle. Each
part stands on its own, and they are ordered the way software actually comes into being.

\begin{description}[leftmargin=0pt, itemsep=4pt, font=\sffamily\bfseries\color{brand}]
  \item[Part I \textperiodcentered{} Product Discovery.] Deciding \emph{what} is worth building,
    before any code. It is the most commonly skipped phase, and the one that stops you building the wrong
    thing well.
  \item[Part II \textperiodcentered{} Project Planning.] Turning a chosen problem into a plan and a
    set of living documents: problem statement, architecture, developer and operational guides.
  \item[Part III \textperiodcentered{} Writing Good Code.] Two chapters: the craft of code that is
    correct, clear, and changeable, and the engineering practices a whole codebase needs, sequenced
    as a maturity model so you know the order in which to adopt them.
  \item[Part IV \textperiodcentered{} Building Real Applications.] A 43-chapter narrative, followed by a project, that takes
    a single Notes API from a command in a terminal to a containerized, HTTPS-served, monitored
    production service, and then on to the distributed systems beyond it. It is divided into eleven
    \emph{stages}, from Stage 1, a program running on one laptop, to Stage 11, systems beyond a single service.
\end{description}

\noindent The closing chapter, \textbf{The Same Process in a Real Team}, replays the whole journey as a
team makes it (who does each step, and what they hand to each other) and lays the same map over
other kinds of software. Every part, and every stage of Part IV, opens by marking where you are on
the map.

\section*{Two ways to read it}

\textbf{To learn,} read the opening chapter first, then Part IV straight through; it builds the
mental model of how a real system works. Then read Part III for the craft, and keep Part II for when you start something of
your own. \textbf{To build something now,} work through Parts I and II to decide what to build and to
frame it, use Part III as your checklist while you write it, and use Part IV as the reference for
taking it all the way to production.

\section*{Exercises, a project, and a companion repository}

Every chapter of Part IV ends with four exercises: \emph{predict} what will happen, \emph{break and
fix} something on purpose, \emph{extend} the code, and \emph{explain} an idea in your own words.
Part IV then closes with a project, \emph{Notes for Many Users}, in five milestones, which you plan
and build yourself. The companion repository holds the Notes API's code, a solution to every
exercise (one file per chapter, in \texttt{exercises/}), and a complete, tested reference solution to
the project, with one Git tag per milestone, so that you can compare your work at each step.
Try each exercise, and each milestone, before you look.

\section*{Conventions used in this book}

Source code appears in shaded listings with a colored rule on the left. Numbered listings carry
a caption:

\begin{codeboxcap}{\textcolor{accent}{Listing 0.1}\enspace A minimal Python program}
\begin{Shaded}
\begin{Highlighting}[]
\KeywordTok{def}\NormalTok{ greet(name: }\BuiltInTok{str}\NormalTok{) }\OperatorTok{{-}\textgreater{}} \BuiltInTok{str}\NormalTok{:}
    \ControlFlowTok{return} \SpecialStringTok{f"Hello, }\SpecialCharTok{\{}\NormalTok{name}\SpecialCharTok{\}}\SpecialStringTok{!"}
\end{Highlighting}
\end{Shaded}
\end{codeboxcap}

Commands you type and the output they produce appear in plain framed blocks:

\begin{outputbox}
\begin{Verbatim}[fontsize=\codesize]
$ python app.py
 * Running on http://127.0.0.1:5000
\end{Verbatim}
\end{outputbox}

Diagrams are drawn in text, so they stay exact and diff-able in the source:

\begin{diagrambox}{250pt}
\begin{Verbatim}[fontsize=\fontsize{9}{8.55}\selectfont]
┌──────────┐     ┌──────────┐     ┌──────────┐
│  Client  │────▶│  Nginx   │────▶│   App    │
└──────────┘     └──────────┘     └──────────┘
\end{Verbatim}
\end{diagrambox}

\begin{notebox}
\textbf{Note:} asides, warnings and version-specific details appear in callouts like this one.
\end{notebox}

Each chapter of Part IV opens with what it will teach you and closes with its key takeaways.

\section*{Versions, tools, and a book that will age}

Software tooling changes faster than books do. The principles here (how a process binds a port,
why a transaction must commit before a response is sent, why a deployment needs a rollback) will
hold for decades. The specifics will not: version numbers, command-line flags, package names,
configuration syntax, default settings, cloud providers' products and prices, and even what counts
as current best practice. Everything version-specific in this book was written and checked against
the versions listed below. When you follow along later, expect some details to differ, and treat
the official documentation of each tool, not this book, as the authority for version-specific
facts. When a command or setting fails, check the tool's current documentation before assuming the
concept is wrong. Prices, free-tier limits, and product names quoted from vendors were correct at
the time of writing and are the details most likely to have changed.

\begin{center}\small
\begin{tabular}{@{}ll@{\quad}ll@{}}
\toprule
Python & 3.11 & Ubuntu Server & 22.04 LTS \\
Flask / Werkzeug & 3.0.2 / 3.0.1 & Docker Engine & 25 \\
Gunicorn & 22.0 & Docker Compose & v2 \\
psycopg2-binary & 2.9.9 & Nginx & 1.18 (Ubuntu 22.04 package) \\
PostgreSQL & 15–16 & Certbot & 1.21 (apt) or current (snap) \\
pytest & 8.1 & GitHub Actions & checkout v4, setup-python v5 \\
bcrypt / requests & 4.1.2 / 2.31.0 & Mailpit & current image \\
prometheus-flask-exporter & 0.23 & openapi-spec-validator & 0.7 \\
\bottomrule
\end{tabular}
\end{center}

Pinned versions in the listings are the ones the examples were written against, not a
recommendation to stay on them. Updating dependencies, and checking them with a vulnerability
scanner as Chapter 17 shows, is part of the job. The companion repository does exactly that: when
the scanner reports a known vulnerability in a version this book pins, the repository moves to a
fixed release, so its versions may be newer than the ones in the listings.
